% C.3aj: full proof independently audited by root, Frobenius and cycles; source checks retained.
\paragraph{C.3aj (Refuted, with fully specified Hodge-generic counterexamples): ordinary higher-cycle restriction in every range $N\ge 2r+2$.}
For every pair of integers $r\ge2$ and $N\ge2r+2$, we specify
an excellent normal strictly henselian local complex ring $R$ of
dimension $N$, with regular puncture $X$, and a nonzero element
$f\in\mathfrak m_R$ whose punctured divisor $Y=X\cap V(f)$ is
regular, such that ordinary algebraically trivial cycles have a
surjection
\[
 \operatorname{coker}\bigl(
 \mathrm{CH}^r(X)_{\mathrm{alg}}\longrightarrow
 \mathrm{CH}^r(Y)_{\mathrm{alg}}\bigr)
 \longrightarrow
 Q_C:=\frac{\operatorname{Pic}^0(C)(\mathbb C)}
                  {\operatorname{Pic}^0(C)(\overline K)}\ne0.
\]
Here $C=V(z_0^4+z_1^4+z_2^4)\subset\mathbb P^2$ and
$\overline K\subset\mathbb C$ is a specified countable algebraically
closed field of definition of the construction. The group $Q_C$ is
uniquely divisible. Thus restriction fails to be surjective even
after tensoring with $\mathbb Q$. Every algebraic equivalence used
below is represented on an actual fixed product with a smooth
connected curve.

\emph{1. The surfaces, coefficient fields, and generic section data.}
Put
\[
 s=r-1,\qquad m=N-2r\ge2.
\]
Choose $s$ smooth quartic K3 surfaces $W_\nu\subset\mathbb P^3$
containing the respective fixed lines
$D_\nu=V(w_{\nu0},w_{\nu1})$, with hyperplane classes $H_\nu$.
The tuple is chosen Hodge-generic in the product of these families,
in the precise sense of the following auxiliary calculation:
\[
 \operatorname{NS}(W_\nu)=\mathbb ZH_\nu\oplus\mathbb ZD_\nu,
 \qquad
 \operatorname{Hg}\left(\bigoplus_{\nu=1}^s T_\nu\right)
       =\prod_{\nu=1}^s\operatorname{SO}(T_\nu),
 \quad
 T_\nu=\operatorname{NS}(W_\nu)_{\mathbb Q}^{\perp}.
\]
The Hodge-generic complex tuple itself is part of the data. We
do not claim that all embeddings of a coefficient field preserve
this Hodge-genericity.

% Auxiliary Hodge input for the higher-range local restriction construction.
% This fragment proves the Hodge assertion only, not the local geometry.
\paragraph{Auxiliary product Hodge calculation for quartics containing lines.}
Let $s\geq1$. We choose smooth complex quartics $W_1,\ldots,W_s$, each
containing a specified line $D_i$, at a jointly Hodge-generic point of
the product of the corresponding quartic families. Write
$H_i=c_1(\mathcal O_{W_i}(1))$ and
$T_i=(\mathbb QH_i+\mathbb QD_i)^\perp\subset H^2(W_i,\mathbb Q)$.
We can and do require
\begin{equation}\label{eq:product-K3-generic-data}
 \begin{gathered}
 \operatorname{NS}(W_i)=\mathbb ZH_i\oplus\mathbb ZD_i,
 \qquad
 \begin{pmatrix}H_i^2&H_iD_i\\H_iD_i&D_i^2\end{pmatrix}
       =\begin{pmatrix}4&1\\1&-2\end{pmatrix},\\
 \operatorname{End}_{\mathrm{Hdg}}(T_i)=\mathbb Q,
 \qquad
 \operatorname{Hom}_{\mathrm{Hdg}}(T_i,T_j)=0\quad(i\ne j).
 \end{gathered}
\end{equation}
Let $C$ be the Fermat plane quartic, set
$S=(\prod_{i=1}^sW_i)\times C$, and consider the indicated Kunneth
summand
\[
 K=\left(\bigotimes_{i=1}^sH^2(W_i,\mathbb Q)\right)
                       \otimes H^1(C,\mathbb Q)
       \subset H^{2s+1}(S,\mathbb Q).
\]
Then its largest rational sub-Hodge structure of level at most one is
\begin{equation}\label{eq:product-K3-level-one}
 K_{\leq1}=
 \left(\bigotimes_{i=1}^s\operatorname{NS}(W_i)_{\mathbb Q}\right)
                       \otimes H^1(C,\mathbb Q).
\end{equation}
In particular, for every smooth projective complex curve $P$ and every
$z\in\mathrm{CH}^{s+1}(S\times P)$, the component of its rational
cycle class in $K\otimes H^1(P,\mathbb Q)$ belongs to
$K_{\leq1}\otimes H^1(P,\mathbb Q)$. The parameter curve is arbitrary;
it is not required to be defined over the field of definition of the
quartics. We assert \eqref{eq:product-K3-level-one} for the specified
Kunneth summand $K$, not for all of $H^{2s+1}(S,\mathbb Q)$.

\emph{Existence of the chosen tuple.}\par
For a fixed line $D_0\subset\mathbb P^3$, let $B$ be the smooth open
subset of $\mathbb P H^0(\mathcal I_{D_0}(4))$. Theorem~1.1 and
Corollary~1.3(c) of \citet[p.~2 of the opened author version]
{frBrevikNollet2011} give $\operatorname{Pic}(W)=\mathbb ZH\oplus
\mathbb ZD_0$ for a very general point of $B$. Their hypotheses hold
because $\mathcal I_{D_0}(3)$ is globally generated and
$H^0(\mathcal I_{D_0}(2))\ne0$. Adjunction gives the intersection matrix
in \eqref{eq:product-K3-generic-data}.

We impose the additional Hodge conditions explicitly. By Local Torelli
and the line-bundle period description in
\citet[Proposition~6.2.8 and Section~6.2.4, pp.~107--108 of the
opened author version]{frHuybrechtsK3Book}, locally marked deformations
preserving $H,D_0$ fill an open subset of the $18$-dimensional period
quadric in $(\mathbb QH+\mathbb QD_0)^\perp_{\mathbb C}$.
They are locally realized in $B$: very ampleness of $H$ is open;
Riemann--Roch and $H\cdot D_0=1$ make the effective class $D_0$ a line;
and a projective coordinate change returns that line to the fixed line.
Thus the product family has independently varying local periods.

For a nonscalar rational endomorphism $a$ of a marked transcendental
space, the Hodge-endomorphism condition forces the period line to be
an eigenline of $a$. The finite union of its eigenspaces cannot contain
the irreducible nondegenerate period quadric, which spans the ambient
projective space. This condition therefore defines a proper locus.
After imposing $\operatorname{End}_{\mathrm{Hdg}}(T_i)=\mathbb Q$,
each $T_i$ is simple: a nontrivial direct summand would give a
nontrivial rational Hodge projector. A nonzero Hodge map between two
$T_i$ would consequently be an isomorphism. For each rational linear
isomorphism $a:T_i\to T_j$, its Hodge locus forces
$a\omega_i\in\mathbb C\omega_j$. The independent variation of the
two period lines makes this a proper locus in the product family.
There are only countably many such rational maps. Applied to the
weight-zero variations $\operatorname{End}(T_i)$ and
$\operatorname{Hom}(T_i,T_j)$, the Hodge-locus theorem
\citep[Theorem~1.1 and Corollary~1.3, p.~2]{frCDK1994}
makes the loci, with their monodromy translates, countable unions of
closed algebraic subsets. Avoiding these loci and the extra Picard
loci gives \eqref{eq:product-K3-generic-data}. This specifies a jointly
Hodge-generic complex point; it does not assert that every embedding
of a coefficientwise algebraically independent parameter field is
Hodge generic.

\emph{Independence of the orthogonal factors.}\par
By \citet[Theorem~2.2.1 and Example~2.2.3, p.~211]{frZarhin1983},
each $T_i$ has rank $20$ and $\operatorname{Hg}(T_i)=\operatorname{SO}(T_i)$.
Let $G=\operatorname{Hg}(\bigoplus_iT_i)$ and
$\mathfrak g=\operatorname{Lie}(G)$. Its projections onto the
absolutely simple Lie algebras
$\mathfrak g_i=\mathfrak{so}(T_i)$ are surjective.
The radical of $\mathfrak g$ maps to a solvable ideal in every
$\mathfrak g_i$, hence is zero. Decompose this semisimple rational
Lie algebra into its simple ideals. Each projection to
$\mathfrak g_i$ is an isomorphism on exactly one of these ideals and
zero on the others. Consequently, if
$\mathfrak g\ne\bigoplus_i\mathfrak g_i$, one simple ideal maps
isomorphically onto two different factors. The projection to that
pair is the graph of a rational Lie-algebra isomorphism
$\varphi:\mathfrak g_i\simeq\mathfrak g_j$.

We record why this would identify the two standard representations;
this is the point at which the rank matters. Over $\overline{\mathbb Q}$,
an automorphism of $\mathfrak{so}_{20}$ integrates to the simply
connected group, and, modulo an inner automorphism, acts by a
symmetry of its $D_{10}$ diagram
\citep[20.33, p.~108, and 21.4, p.~112]{frMilneReductive2018}.
The simple roots are
$\varepsilon_1-\varepsilon_2,\ldots,
\varepsilon_9-\varepsilon_{10},\varepsilon_9+\varepsilon_{10}$
\citep[18.26--18.27, pp.~91--92]{frMilneReductive2018}.
The only nontrivial diagram symmetry swaps its two short terminal
branches and fixes the first node. It therefore fixes the highest
weight $\varepsilon_1$ of the standard representation, whose weights
are $\{\pm\varepsilon_j\}_{j=1}^{10}$. Highest-weight uniqueness
\citep[20.3--20.4, pp.~102--103]{frMilneReductive2018}
shows that twisting the standard representation by any such
automorphism preserves its isomorphism class.

It follows that the intertwiner space
\[
 \{u\in\operatorname{Hom}_{\mathbb Q}(T_i,T_j):
          uX=\varphi(X)u\text{ for every }X\in\mathfrak g_i\}
\]
is nonzero after scalar extension. These are rational linear equations,
so the rational space itself is nonzero. A nonzero intertwiner is an
isomorphism, since both standard representations are absolutely
irreducible. It commutes with $\mathfrak g$ on the pair, hence with
the connected group $G$. Because $T_i,T_j$ have the same weight, it
is a rational Hodge isomorphism, contrary to
\eqref{eq:product-K3-generic-data}. We conclude that
\[
       \operatorname{Hg}\left(\bigoplus_{i=1}^sT_i\right)
                          =\prod_{i=1}^s\operatorname{SO}(T_i).
\]

\emph{The CM factor and the levels of the tensor constituents.}\par
By \citet[Lemma~2, p.~106, and Example~3, pp.~106--107]
{frBeauville2014}, $\operatorname{Jac}(C)$ is isogenous to $E^3$,
where $E=\mathbb C/\mathbb Z[i]$. Put $V=H^1(E,\mathbb Q)$ and
$U=\operatorname{Res}^1_{\mathbb Q(i)/\mathbb Q}\mathbb G_m$.
Then $H^1(C,\mathbb Q)\simeq V^{\oplus3}$ and
$\operatorname{Hg}(V)=U$. The joint Hodge group is a connected
subdirect product of $\prod_i\operatorname{SO}(T_i)$ and $U$.
Its derived group projects onto the full semisimple product and
trivially onto $U$, so it contains the semisimple product. Surjectivity
onto $U$ now gives
\[
 \operatorname{Hg}\left(\bigoplus_iT_i\oplus V\right)
             =\left(\prod_i\operatorname{SO}(T_i)\right)\times U.
\]

For a nonempty subset $I\subset\{1,\ldots,s\}$, put
$T_I=\bigotimes_{i\in I}T_i$ and $t=|I|$.
The representation $T_I$ of the product of the corresponding
orthogonal groups is absolutely irreducible. Over $\mathbb Q(i)$,
$T_I\otimes V$ is the sum of its twists by the two opposite
nontrivial characters of $U$. These are nonisomorphic absolutely
irreducible representations, exchanged by complex conjugation.
Therefore $T_I\otimes V$ is a simple rational Hodge structure.
After the Tate twist supplied by the remaining $s-t$ Neron--Severi
factors, it has a nonzero component of type
\[
                       (s+t+1,s-t),
\]
and hence has level $2t+1\geq3$. Each remaining Neron--Severi
factor has rank two, and $H^1(C)=V^{\oplus3}$; these give
multiplicities, not further simple constituents of smaller level.

Expand $K$ using $H^2(W_i)=\operatorname{NS}(W_i)_{\mathbb Q}\oplus T_i$.
The term with no transcendental factor is precisely the right-hand
side of \eqref{eq:product-K3-level-one}, and has types
$(s+1,s),(s,s+1)$. Every other term is a direct sum of the simple
constituents just described. Semisimplicity of polarizable rational
Hodge structures shows that none of these other terms has a nonzero
sub-Hodge structure of level at most one. This proves
\eqref{eq:product-K3-level-one}.

Finally, the selected Kunneth component of an algebraic cycle class
$\operatorname{cl}(z)$ is a rational Hodge class of type $(s+1,s+1)$.
Poincare duality on $P$ identifies it with a Hodge morphism
\[
                 H^1(P,\mathbb Q)(-s)\longrightarrow K.
\]
Its source has level at most one. Its image therefore lies in
$K_{\leq1}$, proving the stated consequence for every parameter $P$.


Let $K_0$ be generated over $\overline{\mathbb Q}$ by the finitely
many coefficients of these quartic equations. Set
\[
 \begin{gathered}
 A=\prod_{\nu=1}^s W_\nu,\qquad S=A\times C,\\
 M=\left(\boxtimes_{\nu=1}^s\mathcal O_{W_\nu}(1)\right)
                        \boxtimes K_C,\\
 a=c_1(M)=\sum_{\nu=1}^s H_\nu+k_C,
 \qquad k_C=c_1(K_C).
 \end{gathered}
\]
Choose the $(m+1)3\cdot4^s$ complex numbers
\[
 a_{i,\mathbf j,k},\qquad
 0\le i\le m,\quad\mathbf j=(j_1,\ldots,j_s)\in\{0,1,2,3\}^s,
 \quad 0\le k\le2,
\]
algebraically independent over $K_0$, and put
\[
 s_i=\sum_{\mathbf j,k}a_{i,\mathbf j,k}
                 \left(\prod_\nu w_{\nu j_\nu}\right)z_k.
\]
Let $K$ be generated by $K_0$ and
all the $a_{i,\mathbf j,k}$, and let $\overline K$ be its algebraic
closure inside $\mathbb C$. Both fields are countable.

The common zero scheme $B\subset S$ of the $m+1$ sections is
either empty or a smooth complete intersection of codimension
$m+1$. This is a nonempty Zariski-open condition on the tuples:
$M$ is very ample, and successive general hyperplanes give smooth
transverse intersections until their expected dimension is negative,
when their common intersection is empty. The condition is defined
over $K_0$, so algebraic independence ensures it for our tuple.
A second nonempty open condition, needed for the Chern obstruction,
is proved in step 5 below.

Set
\[
 V=\mathbb P^m_x\times S,\quad
 L=\mathcal O_{\mathbb P^m}(1)\boxtimes M,\quad
 \ell=\sum_{i=0}^m x_i s_i,\quad
 H=V(\ell)\subset V,
 \qquad p:H\longrightarrow S.
\]
Over $S\setminus B$, $p$ is the projective bundle of lines in
the kernel of $\mathcal O_S^{m+1}\to M$, with fibre
$\mathbb P^{m-1}$. At a point of $B$, the $ds_i$ are linearly
independent, so $\sum x_i ds_i$ is nonzero for every projective
point $x$. Thus $H$ is smooth everywhere. It is integral: its
projective-bundle open is integral and its complement
$\mathbb P^m\times B$ has codimension $m\ge2$. Its line bundle
of class $h+a$, where $h=c_1(\mathcal O_{\mathbb P^m}(1))|_H$,
is the restriction of the very ample $L$.

\emph{2. The actual local ring and its regular models.}
Let
\[
 B'=\mathbb P^m_x\times A\times\mathbb P^2_z,
 \qquad
 \mathcal N=\mathcal O_{\mathbb P^m}(-1)
       \boxtimes\left(\boxtimes_\nu\mathcal O_{W_\nu}(-1)\right)
       \boxtimes\mathcal O_{\mathbb P^2}(-1),
\]
and write $t$ for the tautological fibre vector. Put
\[
 F=z_0^4+z_1^4+z_2^4,\qquad
 G=z_0^2z_1^2+z_0^2z_2^2+z_1^2z_2^2,
 \qquad Z=V(F+t\ell G)\subset\operatorname{Tot}(\mathcal N).
\]
The expression is a section of the pullback of
$\mathcal O_{\mathbb P^2}(4)$ and is independent of
trivializations. The zero section on $Z$ is exactly $V$, with
normal bundle $L^{-1}$. The differential $dF$ is nonzero there,
so $Z$ is smooth near $V$. On $\ell G\ne0$ it is a graph over
$B'$. The locus $F=\ell G=0$ has codimension two in $B'$
and cannot supply an additional component of this Cartier
hypersurface. Since the hypersurface is Cohen--Macaulay and
generically reduced on its graph open, $Z$ is integral.

Contract the zero section by the Segre coordinates
\[
 y_{i,\mathbf j,k}=t x_i\left(\prod_\nu w_{\nu j_\nu}\right)z_k.
\]
Take the reduced closed affine image $Q$, its finite normalization
$Q^{\mathrm n}$, and the strict henselization of the local ring
at the unique point over the vertex. This defines $R$ without
an unspecified shrinking choice. The proper contraction is an
isomorphism off the zero section. Properness allows a neighbourhood
of the vertex to be chosen on which the model is regular. Its
normalization has only one point over the vertex: the connected
exceptional fibre $V$ maps surjectively onto that finite fibre.
After the strict henselization let $\mathcal Z$ denote this regular
projective model. As in C.3ae, the coordinates $y$ generate
$\mathcal I_V$ on every Segre chart, and every normalization
function in the maximal ideal vanishes on the reduced $V$.
Consequently
\[
 \mathfrak m_R\mathcal O_{\mathcal Z}=\mathcal I_V,
 \qquad X=\mathcal Z\setminus V.
\]
Pointed etale base change with the same residue field preserves
all these statements and the first infinitesimal neighbourhood.
The ring is excellent and normal, and
$\dim R=\dim V+1=m+2r=N$.

The function
\[
 f=\sum_{i,\mathbf j,k}a_{i,\mathbf j,k}y_{i,\mathbf j,k}
\]
pulls back to $t\ell$. Its strict transform is, before the
indicated local base changes, the total space of $L_H^{-1}$
over $H$. Denote its base change by $\mathcal D$. It is regular,
meets $V$ transversely in $H$, and
$Y=\mathcal D\setminus H=X\cap V(f)$ is regular. All the
schemes, line bundles and maps in this step are defined over
$\overline K$ before the final complex and strict-local base changes.

\emph{3. The incidence operator and its boundary on proper families.}
For each $\nu$, put $\tau_\nu=4D_\nu-H_\nu$ and set
$\tau=\prod_\nu\tau_\nu\in\mathrm{CH}^{s}(A)$. Adjunction gives
\[
 \deg(\tau_\nu H_\nu)=0,\qquad
 \deg(\tau_\nu D_\nu)=-9.
\]
For a smooth projective connected parameter $P$, including
a point, define
\[
 \Phi_P:\mathrm{CH}^r(H\times P)\longrightarrow
                  \operatorname{Pic}(C\times P),
 \qquad
 \Phi_P(\alpha)=\operatorname{pr}_{C\times P,*}
                       \bigl(\tau\,p_*(h^{m-1}\alpha)\bigr).
\]
All products with the displayed divisor classes are defined;
no flatness of $p$ is assumed. Since $[H]=h+a$ in
$\mathbb P^m\times S$, the projection formula in that ambient
projective bundle gives
\[
 p_*(h^{m-1})=1,\qquad p_*(h^m)=a,
 \qquad p_*(h^{m+j})=0\quad(j\ge1).
\]
The same identities hold after taking the product with $P$.

Write
\[
 \Psi_P(\gamma)=\Phi_P\bigl((h+a)\gamma\bigr),
 \qquad \gamma\in\mathrm{CH}^{r-1}(H\times P).
\]
We claim that when $P$ is a smooth projective curve the line
bundle $\Psi_P(\gamma)$ has constant fibre class in
$\operatorname{Pic}(C)$. We prove this on its mixed Hodge
component, without asserting that $\gamma$ lifts as a Chow
class to $V\times P$.

Indeed $\dim V=m+2r-1$ and $H$ is a smooth very ample
divisor. Weak Lefschetz gives an isomorphism of rational
Hodge structures
\[
 H^{2r-3}(V,\mathbb Q)\xrightarrow{\ \sim\ }
 H^{2r-3}(H,\mathbb Q).
\]
The exact range also follows directly from C.2b: the complement
is smooth affine, its cohomology vanishes above $\dim V$,
Poincare duality makes its compactly supported cohomology
vanish below $\dim V$, and the pair sequence gives the
displayed isomorphism since $2r-3<\dim V-1$.
Thus the $H^{2r-3}(H)\otimes H^1(P)$ component of the
cycle class of $\gamma$ has a unique ambient lift. Write the
first factor of that lift as
$\zeta=\sum_j h^j\zeta_j$, with
$\zeta_j\in H^{2r-3-2j}(S,\mathbb Q)$. The preceding
pushforwards give on this component exactly
\[
 \Psi_*\zeta
   =\operatorname{pr}_{C,*}
                 \bigl(\tau(2a\zeta_0+a^2\zeta_1)\bigr).
\]
Terms with a negative cohomological degree mean zero.

Every odd cohomology class of $S=A\times C$ has its odd
degree on $C$, because K3 surfaces have no odd cohomology.
The contribution relevant to $H^1(C)$ therefore uses
$H^1(C)$ and total $W_\nu$ degree $2(s-1-j)$ in
$\zeta_j$. A factor of $k_C=c_1(K_C)$ would make the degree
on $C$ exceed two and gives zero. Any $W_\nu$ degree four
already gives zero after multiplication by $\tau_\nu$.
For a remaining term to survive integration over every
$W_\nu$, exactly $j+1$ factors have initial degree zero,
and the required factors of $a^{j+1}$ must insert $H_\nu$
in those positions. Each such term contains
$\deg(\tau_\nu H_\nu)=0$. This proves
$\Psi_*\zeta=0$, for every $\zeta$, not just algebraic ones.
The line bundle $\Psi_P(\gamma)$ has zero mixed Hodge
component. Its classifying map $P\to\operatorname{Pic}(C)$,
after subtraction of one fibre, factors through the Jacobian
of $P$. Its differential is zero, hence the homomorphism of
complex abelian varieties is zero. This proves the claim
integrally on the fibre classes.

\emph{4. Descent of the entire boundary to a countable field.}
It follows that $\Psi=\Psi_{\operatorname{Spec}\mathbb C}$
kills ordinary algebraic-equivalence differences on the
projective variety $H$. Every cycle on $H_{\mathbb C}$ is
ordinarily algebraically equivalent to a cycle defined over
$\overline K$. Here is the exact finite-type justification.
An integral subvariety gives a complex point of a Hilbert
scheme with its fixed Hilbert polynomial. That Hilbert
scheme is projective of finite type over $\overline K$,
with a universal flat closed family
\citep[Theorems 3.1--3.2 and the Hilbert specialization,
pp.~259--260, 265--266]{D_GrothendieckHilbert1961}.
Because $\overline K$ is algebraically closed, the irreducible
component containing the given point is defined over
$\overline K$ and has a $\overline K$-point. The two points
can be joined by a chain of projective integral curves
in that component. Normalize those curves and pull back
the universal flat family. Their fundamental cycles give
actual families on $H$ times smooth projective connected
curves. They prove the asserted algebraic equivalence.
Applying this to the finitely many components of a cycle
and using that $\Psi$ is defined over $\overline K$ yields
\[
 \Psi\bigl(\mathrm{CH}^{r-1}(H)\bigr)
                  \subset\operatorname{Pic}(C)(\overline K).
\]
This statement requires no claim that cycles on a strictly
local puncture itself have a finite-type parameter space.
Only the smooth projective exceptional section $H$ is used.

Chow localization and refined pullback on
$H\subset\mathcal D$ now define
\[
 \rho:\mathrm{CH}^r(Y)\longrightarrow
                  \operatorname{Pic}(C)(\mathbb C)/
                  \operatorname{Pic}(C)(\overline K).
\]
Namely close a cycle on $\mathcal D$, restrict to $H$, and
apply $\Phi$. Changing the closure adds the pushforward of
a class $\gamma\in\mathrm{CH}^{r-1}(H)$; its refined
restriction changes by $-(h+a)\gamma$, whose $\Phi$-image
is zero in the indicated quotient. Localization also proves
independence of rational-equivalence witnesses. A closure
chosen on $\mathcal D\times P$ before taking a fibre gives
the same resulting $\rho$-class after taking that fibre,
because refined pullbacks and proper pushforwards commute.

Every degree class on $C$ is defined over $\overline K$
using a fixed $\overline K$-point. Thus the last quotient is
canonically $Q_C$. It is divisible because the group of
complex points of its Jacobian is divisible. It is
torsion-free: if $n\lambda$ is a $\overline K$-point of the
Jacobian, the finite morphism $[n]$, defined over the
algebraically closed $\overline K$, has all its preimages
over $\overline K$, so $\lambda$ is already such a point.
Consequently $Q_C$ is uniquely divisible.

\emph{5. The three-section obstruction in arbitrary codimension.}
All differentials in this step are relative to $\mathbb C$.
The first neighbourhood of $V$ in $\mathcal Z$ has ideal
$L$ and deformation class
\[
 \kappa=\sum_{i,\mathbf j,k}a_{i,\mathbf j,k}x_i
       \left(\prod_\nu w_{\nu j_\nu}\right)
                   \otimes\delta(z_kG)
       \in H^1(V,T_V\otimes L).
\]
It has only a $C$-tangent component. As calculated in
C.3m and C.3ae, the three classes $\delta(z_kG)$ form a
basis of $H^1(C,\mathcal O_C)$; their Serre-dual basis
is used for $H^0(C,K_C)$.

For each smooth quartic $W_\nu$, the Euler and normal
sequences and Serre duality give
\[
 H^1(W_\nu,\Omega^1_{W_\nu})/\mathbb C H_\nu
       \simeq\mathcal R_{\nu,4}^*,\qquad
 H^1(W_\nu,\Omega^1_{W_\nu}(1))
       \simeq\mathcal R_{\nu,3}^*,
\]
where $\mathcal R_\nu$ is the Jacobian ring of its quartic.
The maps $\mu_{\nu,w}$ of multiplication by a linear form
kill $H_\nu$ and have zero common kernel modulo $H_\nu$:
dually, the maps $w_{\nu j}\mathcal R_{\nu,3}$ span
$\mathcal R_{\nu,4}$. Hence we can choose a linear form
$w_\nu$ with $\mu_{\nu,w_\nu}(\tau_\nu)\ne0$.

Consider the obstruction map on
\[
 \left(\bigotimes_\nu\mathbb C\tau_\nu\right)
                       \otimes H^0(C,K_C).
\]
Its target is
\[
 H^0(\mathbb P^m,\mathcal O(1))\otimes
       \bigotimes_\nu H^1(W_\nu,\Omega^1_{W_\nu}(1)).
\]
For the special triple
$s_i=(\prod_\nu w_\nu)z_i$ for $i=0,1,2$, with all
other $s_i=0$, the three dual basis vectors map to
$x_i\otimes\bigotimes_\nu\mu_{\nu,w_\nu}(\tau_\nu)$.
They are independent. Thus injectivity is a nonempty
rank-three open condition on the full section tuple.
The cohomology maps, divisor classes and minors are
defined over $K_0$ by algebraic de Rham cohomology.
Our algebraically independent coefficients satisfy this
condition and the independent geometric open condition
from step 1. The special tuple is only a rank witness;
its base locus is not asserted to have the required geometry.

Let $P$ be a smooth projective connected curve and
$\Gamma\in\mathrm{CH}^r(X\times P)$. Close it on the
fixed $\mathcal Z\times P$ and let
\[
 b=\sum_{j=0}^{\min(m,r)}h^j b_j,
 \qquad b_j\in\mathrm{CH}^{r-j}(S\times P),
\]
be its refined restriction to $V\times P$.
The necessary Chern obstruction, in its degree-$r$ form, is
\[
 \kappa\mathbin{\lrcorner}\operatorname{cl}_{\mathrm{Hdg}}(b)=0
       \quad\hbox{in }H^{r+1}
       (V\times P,L\otimes\Omega^{r-1}_{V\times P}).
\]
For clarity this uses exactly the exterior-power sequence
\[
 0\longrightarrow L\otimes\Omega^{r-1}_{V\times P}
 \longrightarrow\Omega^r_{\mathcal Z\times P}|_{V\times P}
 \longrightarrow\Omega^r_{V\times P}\longrightarrow0.
\]
Its boundary is contraction with $\kappa$, as is seen by
comparing local splittings of the conormal sequence.
The Hodge Chern character of a vector bundle on
$\mathcal Z\times P$ lifts through the middle term:
the local Atiyah matrices exist before restriction, and
$\operatorname{tr}(\operatorname{At}^r)/r!$ gives the
degree-$r$ component. A rational virtual bundle with
Chow Chern character equal to the chosen closure exists
by the rational Chern-character isomorphism for this
finite-dimensional regular projective scheme over $R$
\citep[Lemma 42.58.1, Tag 0FEY]{cyStacksChernObstruction}.
Naturality under refined restriction identifies its
restricted degree-$r$ character with the class of $b$
\citep[Remark 42.59.9, Tag 0FF8]{cyStacksChernObstruction}.
Equivalently, the finite bundle and cycle data descend
to a smooth pointed etale finite-type complex model.
No sufficiency assertion about deforming cycles is used.

Project the obstruction to the summand
\[
 H^0(\mathbb P^m,\mathcal O(1))\otimes
 \bigotimes_\nu H^1(W_\nu,\Omega^1_{W_\nu}(1))\otimes
 H^1(C,K_C)\otimes H^1(P,\mathcal O_P).
\]
Only the component
\[
 \bigotimes_\nu H^1(W_\nu,\Omega^1_{W_\nu})\otimes
 H^0(C,K_C)\otimes H^1(P,\mathcal O_P)
\]
of $\operatorname{cl}(b_0)$ contributes. The auxiliary
product-Hodge calculation implies that its rational
antecedent has all its $W_\nu$ factors in
$\operatorname{NS}(W_\nu)_{\mathbb Q}$: the tensor
parts containing a transcendental factor have no
level-one Hodge constituent. Components with a
projective-space differential, including every $h^j b_j$
for $j>0$, cannot reach the displayed target because
$\kappa$ contracts only the curve differential.
Every tensor containing an $H_\nu$ factor is killed
by $\mu_{\nu,w}$. The remaining tensor line is
$\bigotimes_\nu\mathbb C\tau_\nu$, and the rank-three
map just proved is injective on its $H^0(C,K_C)$
factor, also after tensoring with $H^1(P,\mathcal O_P)$.
This component is therefore zero.
Its complex-conjugate component with
$H^1(C,\mathcal O_C)\otimes H^0(P,K_P)$ is also zero:
the cycle class, its Kunneth projection and the
Neron--Severi tensor factors are rational. Thus both
halves of the relevant mixed component vanish.

\emph{6. Source-family vanishing and a nonzero target.}
The exact incidence pushforwards give, for the fixed closure
of the preceding source family,
\[
 \Phi_P(b|_{H\times P})
   =\operatorname{pr}_{C\times P,*}\bigl(\tau(b_0+a b_1)\bigr).
\]
The first term has zero mixed component by step 5 and
$\deg(\tau_\nu H_\nu)=0$. The second also has zero mixed
component: the relevant part of $b_1$ has degree
$2r-3=2s-1$ on $S$, with its odd degree on $C$.
The same factor-by-factor degree argument in step 3
forces the single factor of $a$ into a missing $W_\nu$
position and gives $\deg(\tau_\nu H_\nu)=0$; its
$K_C$ component is zero for degree reasons.
Thus this line bundle has constant fibre class on $P$.
Fibre compatibility of $\rho$ proves that every actual
source-family difference maps to zero. A smooth open
parameter curve can be compactified and its cycle closed
on the fixed product, without changing the original
fibres. Consequently
\[
 \rho\bigl(\mathrm{CH}^r(X)_{\mathrm{alg}}|_Y\bigr)=0.
\]

Pull back the divisor classes $D_1,\ldots,D_s$ and a
line bundle from $C$ to $Y$, and multiply their first
Chern classes. This defines
\[
 \theta:\operatorname{Pic}^0(C)(\mathbb C)
          \longrightarrow\mathrm{CH}^r(Y)_{\mathrm{alg}}.
\]
The Poincare line bundle and smooth proper parameter
curves in its connected Jacobian give actual ordinary
families for these cycles. All the pullbacks are of
line bundles, so the possibly nonflat $p$ causes no
undefined operation. Taking the same divisor product
on $\mathcal D$ and using $p_*h^{m-1}=1$ gives
\[
 \rho\theta(\lambda)=(-9)^s\lambda\pmod{
                                      \operatorname{Pic}(C)(\overline K)}.
\]
Since multiplication by the nonzero integer $(-9)^s$
is surjective on $\operatorname{Pic}^0(C)(\mathbb C)$,
this proves the asserted integral cokernel surjection.

Finally choose $\eta\in\overline{\mathbb Q}$ with
$\eta^4=-1$, put $o=(1:\eta:0)$, choose $u$
transcendental over $\overline K$, choose
$v^4=-1-u^4$, and put $q=(u:v:1)\in C(\mathbb C)$.
No nonzero multiple of $[q-o]$ belongs to
$\operatorname{Pic}^0(C)(\overline K)$: finiteness of
$[n]$ would first make its Abel point a $\overline K$-point,
and the Abel closed immersion, defined over that field,
would then make $q$ a $\overline K$-point. This contradicts
the choice of $u$. Thus $Q_C\ne0$, and the target cycle
$\theta([q-o])$ cannot belong to the ordinary source image,
even rationally. The argument proves the stated result
for every $r\ge2$ and $N\ge2r+2$; a finite calculation
is not being used to infer any universal Hodge-genericity.

The retained program
\texttt{code/check\_high\_range\_incidence.py} checks 154
exact pushforward cases, 16 tensor-degree cancellation cases,
eight rank-witness minors and 88 dimension cases. Its runtime
was $0.057480419$ seconds. These checks verify the displayed
coefficient formulas and inequalities in their finite sample;
the arguments above prove the general identities, the actual
family statements and the Hodge-generic existence conditions.
